Sanad hadis merupakan jaringan keilmuan Islam yang terbentuk dari rantai relasi guru--murid antara para perawi hadis lintas generasi. Jaringan ini secara alami dapat dimodelkan sebagai graf berarah, sehingga memungkinkan analisis komputasional untuk mengukur peran dan pengaruh setiap perawi secara kuantitatif. Meskipun sejumlah penelitian terdahulu telah menerapkan pendekatan \textit{Knowledge Graph} dan \textit{Social Network Analysis} (SNA) pada jaringan sanad, hasil analisis yang dihasilkan umumnya bersifat statis dan hanya dapat diakses melalui perangkat lunak desktop khusus. Penelitian ini bertujuan untuk mengembangkan \textit{front-end} berbasis web yang menyajikan \textit{Knowledge Graph} jaringan sanad hadis beserta hasil analisis sentralitas secara interaktif dan dapat diakses publik tanpa instalasi perangkat lunak khusus.

\hspace{1cm} Penelitian dilaksanakan melalui empat tahap utama. Pertama, pra-pemrosesan dataset \textit{Sanadset 650K} yang memuat 650.986 entri hadis dari 926 kitab klasik, meliputi penghapusan baris tanpa sanad, normalisasi penulisan nama perawi, serta resolusi kata kekerabatan. Kedua, pembangunan \textit{Knowledge Graph} dengan mengekstraksi relasi guru--murid dari rantai sanad menjadi tripel terstruktur yang dimodelkan sebagai graf berarah berbobot. Ketiga, analisis jaringan menggunakan lima metrik sentralitas SNA, yaitu \textit{in-degree}, \textit{out-degree}, \textit{eigenvector}, \textit{closeness}, dan \textit{betweenness centrality}, yang dihitung menggunakan pustaka \textit{NetworkX}. Keempat, pembangunan sistem \textit{dashboard} web interaktif menggunakan React dengan basis data Supabase.

\hspace{1cm} Hasil penelitian menunjukkan bahwa dari 487.451 rantai sanad yang telah bersih, terbentuk \textit{Knowledge Graph} berupa graf berarah dengan 177.047 simpul perawi dan 776.798 sisi relasi periwayatan. Analisis SNA mengidentifikasi Abu Hurairah sebagai perawi paling sentral pada empat dari lima metrik, yaitu \textit{in-degree}, \textit{eigenvector}, \textit{closeness}, dan \textit{betweenness centrality}, sementara Sufyan menempati peringkat tertinggi pada \textit{out-degree centrality}. Sistem \textit{dashboard} yang dibangun menyajikan jaringan perawi melalui empat halaman, yaitu Graf Relasi Perawi, Statistik Perawi, Tabel Relasi, dan Tabel Sentralitas, dan telah diterapkan secara daring. Hasil pengujian \textit{black-box} terhadap 14 skenario menunjukkan tingkat keberhasilan 100\%.

\hspace{1cm} Penelitian ini berhasil memenuhi seluruh tujuan yang ditetapkan dan mengisi celah penelitian terdahulu dengan menyajikan \textit{Knowledge Graph} jaringan sanad hadis dalam bentuk \textit{front-end} interaktif yang dapat dieksplorasi secara luas oleh peneliti, mahasiswa, maupun masyarakat umum.

\vspace{0.5cm}
\noindent\textbf{Kata kunci} : \textit{Knowledge Graph}, sanad hadis, \textit{Social Network Analysis}, sentralitas perawi, visualisasi jaringan
